% ==========================================
\section{Sistema Proposto}
% ==========================================

\subsection{Overview}
\textbf{DeQ (Divide et Quantifica)} è un'applicazione web per l'analisi dei mercati e il trading simulato. A differenza delle dashboard classiche a griglia fissa, offre una tela infinita su cui l'utente dispone e collega liberamente dei moduli, i \textit{widget}: azioni, grafici, nodi di calcolo.

Accanto all'analisi c'è un \textbf{Terminale Broker} simulato, dove si comprano e vendono asset con un credito virtuale e una commissione dell'1\% su ogni operazione.

L'app funziona su desktop e mobile ed è pensata per trader, analisti e studenti. Una \textbf{Classifica (Leaderboard)} permette di condividere le proprie schede e confrontarsi con gli altri sul guadagno o la perdita delle posizioni.

\subsection{Requisiti Funzionali}
Il sistema deve permettere all'utente di seguire tutta la propria strategia: idearla, provarla con operazioni simulate e infine condividerla.

Ci sono due tipi di utente: l'\textbf{Utente Standard (Trader)} e l'\textbf{Amministratore}. Chi non è registrato non può usare l'app: per accedere agli strumenti e al listino bisogna registrarsi.

L'\textbf{Utente} deve poter:
\begin{itemize}[itemsep=-2pt]
    \item creare e gestire più \textbf{workspace}, salvati in automatico e quindi disponibili su qualsiasi dispositivo;
    \item mettere \textbf{widget} sulla canvas (sorgenti dati, grafici, nodi di calcolo) e collegarli tra loro;
    \item \textbf{comprare e vendere} asset in modo simulato, su un listino aggiornato;
    \item \textbf{pubblicare le proprie schede in classifica} e lasciare recensioni agli altri utenti.
\end{itemize}

All'\textbf{Amministratore} serve un \textbf{Pannello di Controllo} per gestire gli utenti, aggiungere o togliere titoli dal listino e vedere i ricavi (fittizi) dell'app.

\subsection{Requisiti Non Funzionali}

\subsubsection{Prestazioni e Architettura}
\begin{itemize}
    \item \textbf{Salvataggio automatico:} ogni modifica al workspace viene salvata in background, senza interrompere il lavoro. Lo stato della canvas è convertito in JSON e memorizzato nel database.
    \item \textbf{Cache delle quotazioni:} per non superare i limiti di richieste di Yahoo Finance, il backend conserva per un po' le quotazioni già scaricate e le riusa, riducendo le chiamate esterne.
\end{itemize}

\subsubsection{Sicurezza e Privacy}
\begin{itemize}
    \item \textbf{Sessione singola (JWT):} l'accesso usa token JWT. A ogni nuovo login il token precedente viene invalidato (\texttt{token\_version}), così l'account non resta aperto su più dispositivi.
    \item \textbf{Recupero password sicuro:} il reset avviene tramite un token inviato via email, valido un'ora e utilizzabile una sola volta.
\end{itemize}

\subsubsection{Usabilità e Design}
\begin{itemize}
    \item \textbf{Navigazione fluida:} l'interfaccia gestisce mouse, touch e pennino (\textit{Pointer Events}) per pan, zoom e scorrimento senza scatti.
    \item \textbf{Errori leggibili:} se i dati di mercato non sono disponibili o si inserisce un valore sbagliato, l'app mostra un messaggio chiaro invece di errori tecnici.
\end{itemize}

\subsubsection{Conformità normativa}
\begin{itemize}
    \item \textbf{GDPR (Regolamento UE 2016/679):} eliminando un account vengono cancellati a cascata (\texttt{CASCADE}) tutti i suoi dati: profilo, portafoglio, transazioni e workspace (``diritto all'oblio'').
\end{itemize}

% ==========================================
\section{Macro Pagine e Funzionalità Core}
% ==========================================

\subsection{Autenticazione \& Welcome Page}
È la pagina d'accesso. Per registrarsi bastano nome, username, email e password (con conferma); indirizzo e telefono si aggiungono dopo, dal profilo. C'è anche il recupero della password via email e, dalla stessa schermata, il recupero dell'email collegata all'account.

\subsection{Dashboard (Infinite Workspace)}
Il cuore di DeQ: una tela infinita dove l'utente costruisce la propria analisi, con un'interfaccia scura in stile neon.
\begin{itemize}[itemsep=-2pt]
    \item \textbf{Canvas infinito:} pan, zoom con la rotella e pinch-to-zoom fluidi, su uno sfondo a griglia animata.
    \item \textbf{Schede (tab):} più schede per separare le strategie, con salvataggio automatico in background. Il numero di schede aperte insieme dipende dalla licenza Premium.
    \item \textbf{Annulla (Undo):} l'app tiene la cronologia delle azioni; con \texttt{Ctrl+Z} (o \texttt{Cmd+Z}) si annulla l'ultima operazione e si recupera un widget o un collegamento cancellato per sbaglio.
    \item \textbf{Catalogo widget:} un pannello da cui trascinare i moduli sulla canvas. I widget adattano il testo al livello di zoom e hanno un look \textit{glassmorphism}. Comprende:
    \begin{itemize}
        \item \textbf{Azione:} segue un singolo titolo (quotato o personalizzato) con dati reali.
        \item \textbf{Grafici (ECharts):} candele, linea, area, volumi e torta.
        \item \textbf{Aggregatori:} nodi di calcolo (netto delle posizioni, prezzo medio d'acquisto).
        \item \textbf{Utility:} inventario, orologio e blocchi di testo.
    \end{itemize}
    \item \textbf{Collegamenti:} si clicca per collegare una sorgente a un consumatore; i dati passano lungo linee animate. Un controllo evita i cicli infiniti.
    \item \textbf{Inspector:} un pannello laterale che cambia in base al widget selezionato. Da qui si scelgono i \textit{timeframe} (da 1 giorno a 5 anni), si forzano le connessioni e si cambiano i colori.
\end{itemize}

\subsection{Terminale Broker \& Nastro di mercato}
Il punto in cui l'analisi diventa azione (simulata).
\begin{itemize}[itemsep=-2pt]
    \item \textbf{Terminale e carrello:} listino aggiornato, ricerca dei titoli e checkout per acquisti multipli col credito virtuale. Il saldo si aggiorna da solo, l'operazione viene registrata e si applica l'1\% di commissione.
    \item \textbf{Inventario e storico:} schede per rivedere gli ordini fatti (con mini-grafici) e operare sulle posizioni. Da ogni riga dell'inventario si può comprare, vendere o aggiungere al carrello al volo; se manca il credito appare subito un avviso.
    \item \textbf{Nastro di mercato:} la striscia in basso mostra le variazioni percentuali in tempo reale; cliccando un titolo lo si aggiunge al carrello.
\end{itemize}

\subsection{Classifica (Leaderboard)}
Serve a condividere le strategie e confrontarsi con gli altri. L'utente pubblica una o più schede; la classifica le ordina in base al guadagno o alla perdita delle posizioni, cioè la differenza tra il prezzo attuale e il prezzo medio d'acquisto, senza contare il patrimonio totale. Ogni scheda ha un'etichetta e può essere ritirata quando si vuole; non si può pubblicare due volte la stessa. Si può personalizzare il profilo con un avatar, lasciare recensioni agli altri e comprare in blocco il portafoglio di un altro utente (sempre simulato).

\subsection{Pannello Admin}
Il pannello riservato all'amministratore.
\begin{itemize}[itemsep=-2pt]
    \item \textbf{Gestione piattaforma:} elenca gli utenti registrati e permette di eliminarli, aggiungere o togliere titoli dal listino (collegato a Yahoo Finance) e riempirlo in fretta con titoli scelti a caso da un ampio elenco. Mostra anche lo stato dell'aggiornamento automatico dei prezzi, con conto alla rovescia e possibilità di forzarlo.
    \item \textbf{Finanze:} riepilogo dei ricavi simulati, tra licenze Premium e commissioni sulle operazioni.
\end{itemize}

% ==========================================
\section{Definizioni}
% ==========================================

\begin{table}[H]
\centering
\arrayrulecolor{SlateGray}
\renewcommand{\arraystretch}{1.3}
\begin{tabularx}{\textwidth}{|l|X|}
\hline
\rowcolor[HTML]{3B3766}
\textcolor{white}{\textbf{Termine}} & \textcolor{white}{\textbf{Descrizione}} \\ \hline
\textbf{Utente (Trader)} & Persona registrata che usa la piattaforma per creare strategie e simulare operazioni. \\ \hline
\textbf{Workspace (Canvas)} & La tela infinita, navigabile con pan e zoom, dove l'utente dispone e collega i nodi per comporre la sua dashboard. \\ \hline
\textbf{Widget (Nodo)} & Modulo autonomo (Azione, Grafico, Orologio, Media\ldots) da mettere sulla canvas. Può fare da sorgente o da destinazione dei dati tramite i collegamenti. \\ \hline
\textbf{Terminale Broker} & La parte per comprare e vendere: listino, carrello, inventario e storico, dove si spende il credito virtuale. \\ \hline
\textbf{Leaderboard Entry} & La scheda che un utente rende pubblica per la classifica. Il punteggio è il guadagno o la perdita delle posizioni, cioè la differenza tra prezzo attuale e prezzo medio d'acquisto. \\ \hline
\textbf{Provider Dati} & Servizio esterno (es. Yahoo Finance) da cui il backend prende quotazioni e serie storiche per grafici e listino. \\ \hline
\textbf{Token JWT} & Il token usato per l'accesso. Grazie a un numero di versione garantisce una sola sessione: un nuovo login disconnette gli altri dispositivi. \\ \hline
\textbf{Modale} & Finestra che appare sopra la schermata e blocca il resto finché non ci si interagisce (es. per confermare un acquisto). \\ \hline
\end{tabularx}
\end{table}